\chapter{双勾股定理（The Dual Pythagorean Theorem）}

《正交》（\emph{Orthogonal}）宇宙中的许多物理, 都可以理解为由几何学中相当简单的想法所导出。
尽管由此得到的物理结论时常与我们自身宇宙中的结论大不相同, 几何本身却并无二致。
本章介绍其中一个极其简单的例子, 它却是小说中最重要的科学发现之一背后的基础;
同时, 它也是一个凭其在现实世界中的应用便值得掌握的结果。
我们将先给出这一定理及其几何证明, 随后在矢量空间的语言下, 说明它在对偶空间中的表述。

\section{勾股定理的两种形式}

我们先从一个几乎人人都会一眼解决的问题谈起。

\orthofig{images/01__001.png}{斜向道路的长度}

假设有人告诉我们: 一辆汽车沿一条笔直的道路行驶, 这条路既不完全南北向, 也不完全东西向。
汽车最终停在出发点以北 $3\,\mathrm{km}$、以东 $4\,\mathrm{km}$ 处。那么汽车究竟行驶了多远？

勾股定理（Pythagoras's Theorem）告诉我们, 直角三角形斜边长度的平方等于另外两条直角边长度的平方之和。
于是答案立刻可得: $3$ 与 $4$ 的平方分别为 $9$ 与 $16$, 它们的和为 $25$, 而 $25$ 正是 $5$ 的平方。
所以汽车行驶的总距离为 $5\,\mathrm{km}$。

\orthofig{images/01__002.png}{田地中平行等距的犁沟}

接下来, 考虑一块用平行等距的犁沟耕作过的田地, 这些犁沟既不完全南北向, 也不完全东西向。
从一根钉入地面的木桩出发, 如果我沿正北方向拉出一段绳子, 会发现它恰好穿过 $3$ 条完整的犁沟;
如果我把\textbf{同样长度的一段绳子}改沿正东方向拉出, 会发现它穿过 $4$ 条完整的犁沟。
问题在于: 如果我选择一个方向, 使绳子径直、垂直地横穿犁沟, 那么它会穿过多少条犁沟？

如果你依照上一题类推, 回答"五条", 那么你的答案完全正确——
但如果你以为这不过是勾股定理的又一次直接应用, 那就应当仔细看看插图。
这里\textbf{并不}存在边长为 $3,4,5$ 个单位的直角三角形!

然而这个答案并非误导性的巧合: 若我们把正北、正东方向穿过的犁沟数分别换成任意的 $n$ 与 $e$,
那么同样长度的绳子垂直穿过犁沟的数目 $f$ 将满足
\begin{equation}
f^{2} = e^{2} + n^{2}.
\end{equation}
但若它\textbf{不是}勾股定理的又一个实例, 这个关系究竟从何而来？

\orthofig{images/01__003.png}{田角处的直角三角形}

如果我们考察一个直角三角形, 它的直角顶点在木桩处, 斜边是经过木桩的第一条犁沟,
两条直角边分别沿正北与正东方向, 那么这两条边的长度将分别为 $r/3$ 与 $r/4$,
其中 $r$ 是绳子的长度。

连接该三角形的直角顶点与斜边、并与斜边成直角的那条线段, 其长度为 $r/5$。

\orthofig{images/01__004.png}{带斜边高的直角三角形}

现在不妨转向一般情形: 一个直角三角形, 两条直角边长分别为 $a$ 与 $b$, 斜边长为 $c$,
从直角顶点到斜边的垂直距离为 $d$。

勾股定理给出 $a,b,c$ 之间的关系:
\begin{equation}
c^{2} = a^{2} + b^{2},
\end{equation}
但 $a,b,d$ 之间又有什么关系呢？

虽然 $a,b,c$ 都是三角形的\textbf{边长}, 而 $a,b,d$ 乍看像是毫不相干的三元组, 实际上却并非如此。
它们的共同点在于: 三者都是三角形的\textbf{高}（altitude）——从三角形的某条边到其对顶点的垂直距离。
这为看出它们的关系提供了一条简单的门径: 三角形的面积等于底乘高（相对于该底的高）的一半,
换言之, 等于任意一条边的长度, 乘以从该边到对顶点的垂直距离, 再除以二。

这一事实给出了同一个三角形面积的两种不同表达式, 二者必须相等:
\begin{equation}
\tfrac{1}{2}\, a b = \tfrac{1}{2}\, c d.
\end{equation}
两边同时乘以 $2$, 再平方, 得到
\begin{align}
a^{2} b^{2} &= c^{2} d^{2}, \\
a^{2} b^{2} &= (a^{2} + b^{2})\, d^{2}, \\
\frac{1}{d^{2}} &= \frac{a^{2} + b^{2}}{a^{2} b^{2}},
\end{align}
其中从第一式到第二式用到了勾股定理, 从第二式到第三式则在等式两边同除以 $a^{2} b^{2} d^{2}$。
若将最后一式右端的分母通分消去, 即得

\begin{equation}
\frac{1}{d^{2}} = \frac{1}{a^{2}} + \frac{1}{b^{2}}. \tag{Dual Pythagorean Theorem}
\end{equation}

这一关系适用于任意直角三角形: 其中从两条直角边量得的高为 $a$ 与 $b$, 从斜边量得的高为 $d$。
这正是勾股定理的\textbf{对偶}（dual）形式——它与"平方和"版本 $c^{2}=a^{2}+b^{2}$ 互为对偶:
前者联系三条边长, 后者联系三条高。

现在我们可以\textbf{确切地}看清, 为什么同一段绳子在正北、正东方向穿过的犁沟数 $n,e$,
与垂直穿过的数目 $f$, 满足 $f^{2} = e^{2} + n^{2}$。
木桩处那个三角形对应的距离为
\begin{equation}
a = \frac{r}{e}, \qquad b = \frac{r}{n}, \qquad d = \frac{r}{f},
\end{equation}
代入双勾股定理, 得到
\begin{align}
\frac{1}{d^{2}} &= \frac{1}{a^{2}} + \frac{1}{b^{2}}, \\
\frac{f^{2}}{r^{2}} &= \frac{e^{2}}{r^{2}} + \frac{n^{2}}{r^{2}}, \\
f^{2} &= e^{2} + n^{2}.
\end{align}

\orthofig{images/01__005.png}{地形等高线}

这个结果为什么在物理学中如此重要——无论在《正交》宇宙中, 还是在我们自身的宇宙中？
它所告诉我们的是: 勾股定理的"平方和"公式, 不仅被长度这样的量所遵守,
也被那些"对偶"量所遵守——这些量计数的是\textbf{每}单位长度发生多少次事件, 或\textbf{每}单位长度变化多少。

我们可以把耕过的田地里的犁沟, 换成地形等高线图上的等海拔线。
于是双勾股定理告诉我们: 例如, 若我们所站的地形每向北行进 $1\,\mathrm{km}$ 下降 $30\,\mathrm{m}$,
每向东行进 $1\,\mathrm{km}$ 下降 $40\,\mathrm{m}$,
那么（至少在我们所观察的尺度上, 地形足够接近一个倾斜平面时）
它最陡的坡度, 亦即沿垂直横穿等海拔线方向的梯度, 将是每千米 $50\,\mathrm{m}$。

\orthofig{images/01__006.png}{波前}

然而, 这一结果最重要的应用大概要算波的几何。
若我们把犁沟换成\textbf{波前}（wavefront, 例如海浪的波峰, 或电磁波的波峰）,
那么在给定位移内出现的波前数目就称为\textbf{空间频率}（spatial frequency）,
以区别于计数给定时间内振荡次数的普通\textbf{时间频率}。
双勾股定理告诉我们: 一个波的总空间频率——沿波传播方向测得的每单位距离波前数——
其平方等于沿两个或更多相互垂直的方向做同样测量所得结果的平方之和。

在我们自身的宇宙中, 这已经是一个有用的知识。
而在《正交》的宇宙中, 它彻底革新了人们对物理学的理解。

对于不介意稍高一些数学水平的读者, 我们在下一节提供了关于本主题的进阶材料。

\section{对偶矢量}

\begin{notebox}[进阶内容]
本节面向愿意接受稍高数学水平的读者, 讨论\textbf{对偶矢量}（dual vectors）,
并从对偶空间的角度重新表述双勾股定理。
\end{notebox}

设 $V$ 是一个 $n$ 维矢量空间, $f$ 是从 $V$ 到实数集 $\mathbb{R}$ 的一个线性函数。

$f$ 的\textbf{核}（kernel）, 记作 $\ker f$, 是 $V$ 中使 $f$ 取零的那个子空间。
$\ker f$ 的维数必定至少为 $n-1$（除非 $f$ 是零函数, 此时它恰好等于 $n-1$）。
例如, 若 $n=2$ 且 $f \neq 0$, 则 $f$ 将在 $V$ 中一条过原点的直线上为零;
若 $n=3$ 且 $f \neq 0$, 则 $f$ 将在一个过原点的平面上为零, 依此类推。
这来自线性代数的一个基本定理, 即\textbf{秩-零化度定理}（Rank-Nullity Theorem）。

\orthofig{images/01__007.png}{二维矢量空间上的线性函数}

\orthofig{images/01__008.png}{二维矢量空间上的线性函数}

图中给出一个二维的例子; 这里函数 $f$ 在由矢量 $\bk$ 张成的过原点直线上为零。

若我们选取一个使 $f(\bv)$ 非零的矢量 $\bv$, 那么在 $V$ 中就有一条\textbf{不过原点}、
而是经过 $\bv$ 端点的直线, 其上 $f$ 恒等于 $f(\bv)$。
我们通过给 $\bv$ 加上 $\bk$ 的任意倍数来构造这条直线:
\begin{equation}
f(\bv + s \bk) = f(\bv) + s\, f(\bk) = f(\bv) = 1.
\end{equation}
最后一处等号之所以成立, 只是因为我们恰好选取了满足 $f(\bv)=1$ 的 $\bv$。
不过, $f$ 取常值的所有集合都具有这种形式: 平行于矢量 $\bk$ 的直线。
正如我们在图上把不同的矢量画成长短、方向各异的箭头,
我们也可以把不同的线性函数画成疏密、取向各异的\textbf{线束}（stacks of lines）。

越长的矢量画越大的箭头; 但在这一表示方案里, 一个更大的函数并不是画成间距更大的直线,
而是画成\textbf{排布得更紧密}的直线。
例如, 函数 $g = 2 f$ 的线束中, 每一对直线之间都会多出一条线,
因为当 $f$ 仅为 $1/2$ 时它就已经达到 $1$。

对任意矢量而言, $f$ 的值实际上是在计数该矢量的箭头穿过了线束中的多少个间隙。
例如, $\bw$ 的箭头穿过了 $f$ 的线束中的两个完整间隙, 故 $f(\bw)=2$。
我们还需顾及反向穿过间隙从而给出负值的情形（如 $\bp$）,
以及未穿过整数个间隙的情形（如 $\bq$）;
因此这更像是诉诸几何直观, 而非严格的数学方案。
不过你大概已经能看出, 它如何与主页上的例子相呼应——
在那里, 我们讨论的是在某个方向上移动给定距离时, 穿过了田间的多少条犁沟, 或地图上的多少条等高线。

若 $V$ 的维数大于 $2$, 那么在 $V$ 中我们得到的将不是一束直线, 而是维数为 $n-1$ 的一束平面或超平面。

现在, 考虑从 $V$ 到 $\mathbb{R}$ 的\textbf{所有}线性函数构成的集合。
我们可以把这个集合本身变成为一个 $n$ 维矢量空间, 它称为 $V$ 的\textbf{对偶空间}（dual space）, 记作 $V^{*}$。
$V^{*}$ 中的元素（例如 $f$）称为\textbf{对偶矢量}（dual vectors）。

要使一个集合成为实矢量空间, 我们需要能够把它的元素相加, 并乘以实数。
对 $V$ 上的线性函数集合, 这可以很自然地做到:
若 $f,g \in V^{*}$、$\bv \in V$、$s \in \mathbb{R}$,
我们通过规定它们作为函数取的值, 来定义和 $f+g$ 与标量倍 $s f$:
\begin{align}
(f+g)(\bv) &= f(\bv) + g(\bv), \\
(s f)(\bv) &= s\, f(\bv).
\end{align}

与任何矢量空间一样, 我们可以为 $V^{*}$ 选择一组基, 即 $n$ 个线性函数,
使得 $V^{*}$ 中的任何函数都可以用它们线性表出。
我们把这组基的元素记作 $\{\be^{1}, \be^{2}, \dots, \be^{n}\}$, 并写成
\begin{equation}
f = f_{1}\,\be^{1} + f_{2}\,\be^{2} + \dots + f_{n}\,\be^{n},
\end{equation}
其中 $f_{1}$ 等是 $f$ 关于这组基的\textbf{分量}。
在本注记中, 我们通常用上标标记矢量的分量, 却用\textbf{下标}标记对偶矢量的分量;
并且, 我们通常用下标标记基中的矢量, 却用\textbf{上标}标记基中的对偶矢量。

\orthofig{images/01__009.png}{对偶基}

若 $V$ 有一组特定的基 $\{\be_{1}, \be_{2}, \dots, \be_{n}\}$,
那么当 $V^{*}$ 的一组基 $\{\be^{1}, \be^{2}, \dots, \be^{n}\}$ 满足下述条件时,
我们说它\textbf{对偶于}（dual to）$V$ 的这组基:
\begin{equation}
\be^{i}(\be_{j}) = \delta^{i}_{j},
\end{equation}
其中 $\delta^{i}_{j}$ 称为\textbf{克罗内克 $\delta$ 符号}（Kronecker delta symbol）,
当 $i=j$ 时为 $1$, 当 $i \neq j$ 时为 $0$。
从几何上看, 这意味着: 对某个特定的 $i$, 函数 $\be^{i}$ 在 $V$ 中表现为一束直线或平面,
使得矢量 $\be_{i}$ 恰好穿过相邻两线之间的一个间隙,
而所有其他基矢量 $\be_{j}$ 都落在过原点的那条直线或平面之内, 完全不穿过该线束。
例如, 在左图中, 矢量 $\be_{2}$ 落在直线 $\be^{1}=0$ 之内, 而矢量 $\be_{1}$ 落在直线 $\be^{2}=0$ 之内。

给定 $V$ 的基 $\{\be_{1}, \be_{2}, \dots, \be_{n}\}$, 我们只需把原基中的每个矢量 $\be_{i}$ 代入 $f$,
即可求得 $V^{*}$ 中任意线性函数 $f$ 关于对偶基的分量。
因为两组基互为对偶, 只有特定的那个基函数 $\be^{i}$ 在 $\be_{i}$ 处非零,
于是求和中只剩下乘以它的那一项分量 $f_{i}$ 存活下来:
\begin{equation}
f(\be_{i}) = (f_{1}\be^{1} + f_{2}\be^{2} + \dots + f_{n}\be^{n})(\be_{i}) = f_{i}.
\end{equation}

若矢量 $\bv \in V$ 关于 $V$ 某组选定基的分量为 $v^{i}$,
而对偶矢量 $f \in V^{*}$ 关于对偶基的分量为 $f_{i}$, 则
\begin{equation}
f(\bv) = f(v^{i} \be_{i}) = f_{i}\, v^{i},
\end{equation}
其中我们用了\textbf{爱因斯坦求和约定}（Einstein summation convention）来简写对重复指标的求和
（即 $v^{i}\be_{i} = v^{1}\be_{1} + v^{2}\be_{2} + \dots + v^{n}\be_{n}$）。

假设 $V$ 上有一个点积。那么对 $V$ 中任意矢量 $\bw$, 我们可以定义一个 $V^{*}$ 中的线性函数 $f_{\bw}$:
\begin{equation}
f_{\bw}(\bv) = \bv \cdot \bw.
\end{equation}
反过来, 给定 $V^{*}$ 中的任意线性函数 $f$, 我们可以找到这样的矢量 $\bw_{f}$:
使得对每个 $\bv \in V$ 都有 $f(\bv) = \bv \cdot \bw_{f}$。
为此, 我们先选取 $V$ 的一组正交归一基 $\{\be_{1}, \be_{2}, \dots, \be_{n}\}$, 然后令
\begin{equation}
\bw_{f} = f_{i}\, \be_{i},
\end{equation}
这里仍使用爱因斯坦求和约定, 而分量 $f_{i}$ 是相对于与 $V$ 的正交归一基对偶的那组 $V^{*}$ 基取的。于是
\begin{align}
\bv \cdot \bw_{f}
&= \bv \cdot [f_{i} \be_{i}] \\
&= f_{i}\,[\bv \cdot \be_{i}] \\
&= f_{i}\, v^{i} = f(\bv).
\end{align}
因此, 我们可以把 $V$ 中的任一元素与 $V^{*}$ 中的唯一元素对应起来, 反之亦然。

我们用这种方式与 $f$ 相联系的矢量 $f_{i} \be_{i}$,
将正交于任何落在 $f$ 的核中的矢量 $\bk$,
因为 $\bk$ 与 $f_{i} \be_{i}$ 的点积恰好就是 $f(\bk)=0$。
由于我们用来可视化 $f$ 的线束中所有的直线（或平面等）都平行于过原点的那条核直线,
矢量 $f_{i} \be_{i}$ 将垂直于整个线束。

\begin{notebox}[补充: 与洛伦兹情形的对比]
上述 $V$ 与 $V^{*}$ 之间的对应, 本身是在一个\textbf{正定}（黎曼型, 度规全为正是 $(+,+,+,+)$）点积下得到的:
在正交归一基中, 这种对应使矢量与对应余矢量的分量\textbf{逐一相等}, 即 $v_{i} = v^{i}$,
因而两者可以被"等同"起来。
而在本宇宙常用的\textbf{洛伦兹}（Lorentzian）情形中, 度规取 $(-,+,+,+)$,
矢量与余矢量之间只相差一个\textbf{时间分量上的符号}: 若某矢量的分量为 $(v^{x}, v^{y}, v^{z}, v^{t})$,
则由度规降指标所得的对应余矢量分量为 $v_{a} = g_{ab}\, v^{b}$, 即
\[
v_{x} = v^{x}, \qquad v_{y} = v^{y}, \qquad v_{z} = v^{z}, \qquad v_{t} = -\,v^{t}.
\]
因此, 在洛伦兹情形中, 对偶矢量与原始矢量不能再像黎曼情形那样被简单地等同起来,
两者的时间分量相差一个符号。
\end{notebox}

我们可以给 $V^{*}$ 赋予一个点积, 办法是规定:
与 $V$ 的任意正交归一基对偶的任意一组基, 本身也是正交归一的。
这让我们得以经由其平方来定义对偶矢量 $f$ 的"长度"或大小:
\begin{equation}
|f|^{2} = f \cdot f
= (f_{1}\be^{1} + f_{2}\be^{2} + \dots + f_{n}\be^{n})
  \cdot (f_{1}\be^{1} + f_{2}\be^{2} + \dots + f_{n}\be^{n})
= (f_{1})^{2} + (f_{2})^{2} + \dots + (f_{n})^{2}.
\end{equation}
此处 $f_{i}$ 是相对于 $V^{*}$ 的基 $\{\be^{1}, \be^{2}, \dots, \be^{n}\}$ 取的分量,
而该基对偶于 $V$ 的一组正交归一基。

\orthofig{images/01__006.png}{波前}

按照我们对 $f$ 的几何解释, 这里每个分量 $f_{i} = f(\be_{i})$ 都是在计数基矢量 $\be_{i}$ 穿过了多少个"线束中的间隙"。
所以它们就像右图中所画的: 当我们沿 $x$ 轴与 $y$ 轴各移动一米时, 穿过了多少个波的周期。

但现在假设我们径直地、\textbf{垂直}地穿过线束,
用的单位矢量沿着 $f_{i} \be_{i}$ 的方向——确切地说, 是 $f_{i} \be_{i} / |f_{i} \be_{i}|$——
而我们知道它正交于整个线束。这样一次垂直穿越所得到的计数为
\begin{align}
f\!\left(\frac{f_{i} \be_{i}}{|f_{i} \be_{i}|}\right)
&= \frac{f_{i}\, f(\be_{i})}{|f_{i} \be_{i}|} \\
&= \frac{f_{i}\, f_{i}}{|f_{i} \be_{i}|} \\
&= \frac{|f|^{2}}{|f|} \\
&= |f|.
\end{align}

这正是双勾股定理的又一种表述!
用一个单位矢量径直穿过线束所得到的计数为 $|f|$,
而我们已经看到, 它的平方恰好等于分别沿 $n$ 个相互正交方向的单位矢量穿过该线束所得计数的平方之和。
换言之, 在对偶空间中, 那个"平方和"形式（勾股定理的原始版本）再次浮现——
这正是"双勾股"之名所传达的对偶关系: 矢量空间中的长度, 与对偶空间中的计数, 服从彼此对偶的两条勾股关系。
